\section{Conclusions and Priority Experiments}

The available evidence supports the following conclusions. The low-temperature--high-temperature polymorphism of BaZn$_2$Si$_2$O$_7$ is governed by the ZnO$_4$ chains and the corner-sharing SiO$_4$ network; Ba\slash Sr, Zn\slash Mg\slash Co\slash Ni\slash Cu, and Si\slash Ge substitution can tune the transition temperature and thermal expansion systematically; and the solid-state, sol--gel, and glass-crystallization routes each have well-defined kinetic windows but all demand that phase assemblage and microstructure be reported together \cite{lin1999phase,kerstan2012thermal,thieme2015ba1,thieme2016new,bocker2012in}. For Ba$_5$Y$_{12}$Zn[O(SiO$_4$)]$_8$ there are only clues from compositionally similar systems and no verifiable direct structural or synthesis report, so ``structural analogue'' should be treated as a hypothesis awaiting verification.

The highest-priority follow-up experiments, in order, are: (i) determine the ternary BaO--ZnO--SiO$_2$ phase diagram and build a temperature-dependent thermodynamic model; (ii) search a composition gradient for the single-phase window of Ba$_5$Y$_{12}$Zn[O(SiO$_4$)]$_8$ and follow its transitions by HT\nobreakdash-XRD; and (iii) refine the Y\slash Zn coordination and oxygen stoichiometry of candidate samples. Only once these three steps are complete can one judge at what level the structure--processing trends of BaZn$_2$Si$_2$O$_7$ carry over to the Y-rich framework.
